\documentclass[preprint,12pt]{elsarticle}

\usepackage{graphicx}
\usepackage{subcaption}
\usepackage{adjustbox}
\usepackage{amssymb}
\usepackage{amsmath}

\journal{Journal Name}

\begin{document}

\begin{frontmatter}

%% Title, authors and addresses

\title{Description of GPU Navier-Stokes Solver}

\author{E.A.J.F.\ Peters}

\address{}

\begin{abstract}
This paper presents a GPU-accelerated Incompressible Navier-Stokes solver utilizing a staggered MAC grid and a Robust Scaled Ghost-Cell Immersed Boundary Method (IBM). The system is designed to handle complex porous geometries defined by a Signed Distance Function (SDF). High performance is achieved through a specialized matrix-free implementation on the GPU.
\end{abstract}

\begin{keyword}
GPU Computing \sep Navier-Stokes \sep Immersed Boundary Method \sep Signed Distance Function \sep Staggered Grid \sep Porous Media
\end{keyword}

\end{frontmatter}

%%
%% Start line numbering here if you want
%%
%\linenumbers
\section{Navier-Stokes Projection Scheme}

This work addresses the numerical solution of the unsteady Navier-Stokes equations subject to the incompressibility constraint. The solver is designed for high-performance execution on GPUs, employing a standard Marker-and-Cell (MAC) staggered grid arrangement. The primary variables are defined as follows: the velocity components $u$, $v$, and $w$ are located at the face centers $(i+1/2, j, k)$, $(i, j+1/2, k)$, and $(i, j, k+1/2)$ respectively, while the pressure $p$ is defined at the cell centers $(i, j, k)$.

\subsection{Numerical Discretization}
The momentum equation is discretized using a finite difference approach. For the advection terms, we employ a Total Variation Diminishing (TVD) scheme with the Koren limiter. To handle complex geometries and ensure high performance on the GPU, we introduce a \textbf{Symmetric Branchless} formulation for the ghost-cell IBM using a \textbf{Robust Scaled Scheme}.

We formulate the solution procedure as a Newton-Raphson iteration to find the root of the residual vector $\mathbf{f}(\mathbf{u}) = 0$. Using the robust scaling approach, the entire discrete momentum equation is multiplied by the regularization factor $\mathcal{D}_{rescale}$:
\begin{multline}
    \mathbf{f} = \mathcal{D}_{rescale} \left[ \frac{\rho}{\Delta t} (\mathbf{u} - \mathbf{u}_\mathrm{old}) + \theta \, \nabla p + (1-\theta) \, \nabla p_\mathrm{old} \right] \\
    + \theta \left( \rho \, \mathbf{c}_{scaled}(\mathbf{u}) - \mu \, \mathbf{L}_{u,scaled} \, \mathbf{u} \right)
    \\+ (1-\theta) \left( \rho \, \mathbf{c}_{scaled}(\mathbf{u}_\mathrm{old}) - \mu \, \mathbf{L}_{u,scaled} \, \mathbf{u}_\mathrm{old} \right)
\end{multline}
where the scaled spatial operators are defined such that their coefficients include the regularization factor. A key property is that the ratio $\mathcal{D}_{rescale}/\mathcal{D}$ appearing in the stencil weights is strictly bounded ($\le 1$) and computed directly, avoiding all singularities.

To solve this system, we invert the Jacobian $\mathbf{J}$. To ensure diagonal dominance and solver robustness, we approximate the advection part using the linear First-Order Upwind operator:
\begin{equation}
    \mathbf{J} \approx \frac{\rho \mathcal{D}_{rescale}}{\Delta t} \mathbf{I} + \theta (\rho \, \mathbf{C}_{u,scaled} - \mu \, \mathbf{L}_{u,scaled})
\end{equation}

\subsection{Non-Linear Advection Implementation}
The non-linear advection operator $\mathbf{c}_{scaled}(\mathbf{u})$ is decomposed into a robust linear part, suitable for the implicit Jacobian, and a deferred correction source term:
\begin{equation}
    \mathbf{c}_{scaled}(\mathbf{u}) = \mathbf{C}_{u,scaled} \, \mathbf{u} + \mathbf{c}_{corr, scaled}(\mathbf{u})
\end{equation}
The linear operator $\mathbf{C}_{u,scaled}$ corresponds to the First-Order Upwind (FOU) scheme (scaled by $\mathcal{D}_{rescale}$). The correction term encapsulates the difference between the high-order TVD scheme and the FOU operator:
\begin{equation}
    \mathbf{c}_{corr, scaled} = -\mathcal{D}_{rescale} \, \nabla \cdot  \mathbf{F}_{corr}
\end{equation}
where $\mathbf{F}_{corr}$ represents the unscaled difference between the TVD and FOU fluxes:
\begin{equation}
   F_{corr, face} = \frac{1}{2} \Psi(r) (F_{CDS} - F_{FOU})
\end{equation}
The Koren limiter function is given by $\Psi(r) = \max \left( 0, \min \left( 2r, (1+2r)/3, 2 \right) \right)$. The calculation of this correction term requires careful handling of the Immersed Boundary to avoid singularities:
\begin{itemize}
    \item \textbf{Far-Point Solid}: If the upwind node required for the gradient ratio $r$ lies inside the solid, the limiter logic is bypassed. We switch to Central Differencing by enforcing $\Psi=2$, which effectively removes the dependency on the invalid far-upwind node.
    \item \textbf{Near-Point Solid (Distributive Scaling)}: If the stencil involves a ghost-cell near the interface, the flux $F_{corr}$ naturally contains a singularity ($\sim 1/\mathcal{D}_{stencil}$). We handle this by distributing the global scaling factor $\mathcal{D}_{rescale}$ into the divergence summation. For each face $f$, the contribution is computed as:
    \begin{equation}
        \text{Contribution}_f = \text{clamp}\left( \frac{\mathcal{D}_{rescale}}{\mathcal{D}_{stencil}} \right) \cdot \mathcal{F}_{num}
    \end{equation}
    where $\mathcal{F}_{num}$ is the numerator of the interpolation polynomial (bounded) and the ratio of scaling factors is explicitly clamped ($\le 1$). This ensures that the correction term remains well-behaved throughout the domain without ever forming the singular fraction explicitly.
\end{itemize}

For time integration, we introduce a parameter $\theta \in [0.5, 1]$. Setting $\theta=1$ recovers the Fully Implicit scheme, known for its robustness, while $\theta=0.5$ yields the second-order Crank-Nicolson method.

\subsection{Pressure Solver}
The pressure Poisson equation is solved using a Red-Black Gauss-Seidel (RB-GS) smoother. This iterative solver is particularly well-suited for GPU architectures due to its inherent parallelism and "matrix-free" implementation, avoiding the explicit storage of the large system matrix.

\subsection{Projection Step}
To decouple the system, we employ a projection-like update. First, we compute an intermediate velocity increment, $\delta\mathbf{u}^*$, ignoring the pressure update. Subsequently, we introduce an auxiliary scalar field, $\phi$, such that adding its gradient projects the velocity onto the divergence-free subspace:
\begin{align}
    \left( \frac{\rho \mathcal{D}_{rescale}}{\Delta t} \mathbf{I} + \theta (\rho \, \mathbf{C}_{u,scaled} - \mu \, \mathbf{L}_{u,scaled}) \right) \delta \mathbf{u}^* &= - \mathbf{f}^{(k)} \\
    \mathbf{D} \mathbf{G} \, \phi &= \mathbf{D} \, (\mathbf{u}^{(k)} + \delta \mathbf{u}^*) \label{eq:pressure-poisson}\\
    \delta \mathbf{u}^{(k+1)} &= \delta \mathbf{u}^* - \mathbf{G} \, \phi
\end{align}
Here, the second line is a Poisson-type equation for $\phi$, derived by requiring that $\mathbf{D} \delta \mathbf{u}^{(k+1)} = -\mathbf{D} \mathbf{u}^{(k)}$.

Substituting the corrected velocity $\delta \mathbf{u}^{(k+1)}$ back into the linearized momentum equation roughly implies $\delta \mathbf{u} \approx - (\mathbf{J})^{-1} (\mathcal{D}_{rescale} \mathbf{G} \delta p)$. In the pressure correction step, the global scaling factor $\mathcal{D}_{rescale}$ effectively cancels out with the scaled mass matrix term $\frac{\rho \mathcal{D}_{rescale}}{\Delta t}$, preserving the standard structure of the pressure Poisson equation. The simplified update becomes:
\begin{equation}
    \delta \tilde{p}^{(k+1)} = \left( \frac{\rho}{\theta \, \Delta t} \mathbf{I} + \rho \, \mathbf{C}_\phi - \mu \, \mathbf{L}_\phi \right) \, \phi
    \label{eq:pressure_from_phi}
\end{equation}
The velocity and pressure updates are iterated until convergence.

\section{Immersed Boundary Method Treatment}

\subsection{Geometric Representation and Grid Classification}
The solid geometry is represented implicitly using a global Signed Distance Function (SDF), denoted as $\psi(\mathbf{x})$. The SDF indicates the shortest distance to the solid surface, with the convention that $\psi < 0$ inside the solid and $\psi > 0$ in the fluid. 

The solver employs a staggered MAC grid. The primary SDF values are defined at cell centers. Values at staggered velocity locations (faces) are obtained via linear interpolation. The \texttt{classify\_cells\_kernel} is executed to determine the status of each control volume. A velocity component is classified as \textit{fluid} if its corresponding (interpolated) SDF value is positive ($\psi > 0$), and \textit{solid} otherwise ($\psi < 0$). For cell volumes and face areas, the SDF is further utilized to compute fluid volume fractions and area fractions, respectively.


\subsection{Directional Ghost-Cell Methodology}
To impose no-slip boundary conditions with second-order spatial accuracy, we employ a Directional Ghost-Cell Immersed Boundary Method (GC-IBM). When a fluid velocity variable $u_c$ (where $u$ represents any velocity component $u, v, \text{or } w$) has a neighbor classified as solid, it is treated as a "ghost" value. 

We construct a quadratic polynomial fit that satisfies the fluid neighbor value $u_{nb}$, the central value $u_c$, and the physical boundary condition $u_{bc}$ at the surface location $\xi$. The ghost value $u_\mathrm{ghost}$ is then expressed as a function of these known quantities:
\begin{equation}
    u_\mathrm{ghost} \approx w_{nb} \, u_{nb} + w_c \, u_c + w_{bc} \, u_{bc} \label{eq:ghost}
\end{equation}
Standard calculation of these weights $w$ introduces singularities when the boundary distance $\xi \to 0$.

\subsection{Robust Scaled Formulation}

To ensure numerical stability, we implement a \textit{Robust Scaled Scheme}. We define polynomial numerators ($\mathcal{N}$) and a common denominator ($\mathcal{D}$) such that $w = \mathcal{N}/\mathcal{D}$. The discrete governing equation is formally multiplied by $\mathcal{D}$ to eliminate rational singularities.

The geometric functions for the directional neighbor ($\mathcal{N}_{nb}$), center ($\mathcal{N}_c$), and inhomogeneous boundary term ($\mathcal{N}_{bc}$) are defined in Table~\ref{tab:ibm_coeffs}. Note that $\mathcal{N}_{nb}$ corresponds to the fluid neighbor opposite to the ghost cell (e.g., if east is ghost, $\mathcal{N}_{nb}$ applies to West).

\begin{table}[h!]
    \centering
    \renewcommand{\arraystretch}{1.5}
    \begin{tabular}{|l|c|c|c|c|}
        \hline
        \textbf{Scheme} & $\mathcal{D}(\xi)$ & $\mathcal{N}_{nb,d}(\xi)$ &
        $\mathcal{N}_c(\xi)$ &
        $\mathcal{N}_{bc}(\xi)$ \\
        \hline
        Point-value & $\xi(1+\xi)$ & $\xi(1-\xi)$ & $2(\xi^2 - 1)$ & $2$ \\
        \hline
        Cell-avg & $\xi(1+\xi) - \frac{1}{12}$ & $\xi(1-\xi) + \frac{1}{12}$ & $2(\xi^2 - 1) - \frac{1}{6}$ & $2$ \\
        \hline
    \end{tabular}
    \caption{Geometric scaling factors for the robust GC-IBM. $\mathcal{N}_{nb}$ applies to the fluid neighbor used in the interpolation stencil.}
    \label{tab:ibm_coeffs}
\end{table}

\subsection{Modified Matrix Equation}
In the linear solver, we construct the modified system matrix $\mathbf{A}_{IBM}$ and an inhomogeneous term vector $\mathbf{a}_{inhom}$. The original equation $\mathbf{A} \mathbf{x} = \mathbf{b}$ is transformed to:
\begin{equation}
    \mathbf{A}_{IBM} \mathbf{x} + \mathbf{a}_{inhom} = \mathcal{D}_{rescale} \mathbf{b}
\end{equation}
The modified row coefficients $a'$ and the inhomogeneous contribution $a_{inhom}$ for a cell $c$ are computed as follows:
\begin{equation}
    \begin{aligned}
    a'_c &= \mathcal{D}_{rescale} \, a_c + \sum_d \mathcal{N}_{c,d} \, \mathcal{R}_d \, a_{ghost,d}\\
    a'_{nb,d} &= \mathcal{D}_{rescale} \, a_{nb,d} + \sum_d \mathcal{N}_{nb,d}  \, \mathcal{R}_d \, a_{ghost,d} \\
    a'_{nb,\mathrm{other}} &= \mathcal{D}_{rescale} \, a_{nb,\mathrm{other}} \quad \text{(other directions)} \\
    a'_{ghost,d} &= 0 \\
    a_{inhom,c} &= \sum_d \mathcal{N}_{bc,d} \, \mathcal{R}_d \, u_{bc,d} \, a_{ghost,d}
    \end{aligned}
\end{equation}
where $a_{ghost,d}$ is the original stencil coefficient for the neighbor in direction $d$ which is now classified as a ghost cell. For stationary objects with no-slip conditions ($u_{bc}=0$), $a_{inhom}$ vanishes.
When there are only solid neighbors in one direction, the sum over $d$ consist only of that direction, $\mathcal{D}_{rescale} = \mathcal{D}$, as introduced above, and $\mathcal{R}_d =1$.

\subsubsection{Double-Sided Boundary (Sandwiched Cell)}
In packed geometries, a fluid cell $c$ may be bounded by solid surfaces on both sides, e.g., West at distance $\xi_-$ and East at distance $\xi_+$. In this scenario, both neighbors $u_{nb,-}$ and $u_{nb,+}$ are ghost cells.

The system row for cell $c$ is modified to eliminate both neighbors:
\begin{equation}
    a'_c u_c + a_{inhom,c} = \mathcal{D} b_c
\end{equation}
The modified coefficient and inhomogeneous term are derived as:
\begin{align}
    a'_c &= \mathcal{D} \, a_c +  \mathcal{N}_{c,-} \, a_{ghost,-} + \mathcal{N}_{c,+} \, a_{ghost,+} \\
    a_{inhom,c} &= \mathcal{N}_{bc,-,-}  \, u_{bc,-} \,  a_{ghost,-} + \mathcal{N}_{bc,-,+}  \, u_{bc,-} \, a_{ghost,+} \nonumber \\
    & \quad
        + \mathcal{N}_{bc,+,-}  \, u_{bc,+} \, a_{ghost,-}
     + \mathcal{N}_{bc,+,+}  \, u_{bc,+} \,  a_{ghost,+}
\end{align}
The polynomials for the double-sided case are provided in Table \ref{tab:ibm_double}. 

\begin{table}[h!]
    \centering
    \renewcommand{\arraystretch}{2.0}
    \resizebox{\textwidth}{!}{%
    \begin{tabular}{|l|c|c|}
        \hline
        \textbf{Factor} & \textbf{Point-value} & \textbf{Cell-average} \\
        \hline
        $\mathcal{D}$ & $\xi_- \xi_+$ & $(\xi_- \xi_+ - \frac{1}{12})$ \\
        \hline
        $\mathcal{N}_{c,+}$& $(\xi_- + 1)(\xi_+ - 1)$ & $(\xi_- + 1)(\xi_+ - 1) - \frac{1}{12}$ \\
        \hline
        $\mathcal{N}_{bc,+,+}$& $\frac{\xi_-}{\xi_- + \xi_+} (1 + \xi_-)$ & $\frac{\xi_-}{\xi_- + \xi_+} (1 + \xi_-) - \frac{1}{12}$ \\
        \hline
        $\mathcal{N}_{bc,-,+}$ & $\frac{\xi_+}{\xi_- + \xi_+} (1 - \xi_+)$ & $\frac{\xi_+}{\xi_- + \xi_+} (1 - \xi_+) + \frac{1}{12}$ \\
        \hline
    \end{tabular}%
    }
    \caption{Geometric polynomials for the double-sided (sandwiched) case. The cases for $- \leftrightarrow +$ can be obtained by interchanging $\xi_- \leftrightarrow \xi_+$.}
    \label{tab:ibm_double}
\end{table}

\subsection{Multi-Directional IBM}

For boundaries that are not grid-aligned, the IBM stencil often requires corrections in multiple spatial directions. In such cases, we scale by the factor:
\begin{equation}
  \mathcal{D}_\mathrm{rescale} = \mathcal{D}_{d_{min}}, \text{ where } d_{min} = \mathrm{argmin}_{d} (|\mathcal{D}_d|) \label{eq:D_rescale}
\end{equation}
We define the scaling ratio for direction $d$ as $\mathcal{R}_d = \mathcal{D}_{rescale}/\mathcal{D}_d$.
The summation runs over all directions where the IBM is applied. For directions where $d \neq d_{min}$, we have $|\mathcal{R}_d| \le 1$. A singularity ($0/0$) may seemingly arise if $\mathcal{D}_\mathrm{rescale} \approx 0$ and $\mathcal{D}_{d} \approx 0$. This occurs when the IBM interaction point is at (for point values) or very close to (for cell averages) the cell center. In such instances, we define $\mathcal{R}_d = 1$. Since the interaction points in all directions are spatially proximate, the boundary values $u_{bc,d}$ are similar, and the boundary condition remains well-posed.

\subsection{IBM for Pressure Computations}

In incompressible systems, pressure acts as a Lagrange multiplier enforcing the incompressibility constraint. Pressure values are required in cells where any face contains a ``fluid'' velocity, as the pressure gradient must be defined there. The pressures are updated via the projection scalar $\phi$, necessitating the computation of the velocity divergence.

The divergence is calculated based on superficial velocities (or fluxes), $\tilde{\mathbf{u}}$, which are defined as the intrinsic fluid velocity $\mathbf{u}$ weighted by the fluid area fraction $\alpha_f$ of the corresponding face:
\begin{equation}
    \tilde{u} = \alpha_f u
\end{equation}
The area fraction $\alpha_f$ for a face (e.g., with normal in the $x$-direction) is approximated by assuming a planar interface:
\begin{equation}
    \alpha_f \approx \text{clamp}\left( 0.5 + \frac{\psi}{ |n_y|\Delta y + |n_z|\Delta z } , 0, 1 \right)
\end{equation}
where $\psi$ is the Signed Distance Function (SDF) evaluated at the face center, and $\mathbf{n}$ is the interface normal vector. This formula represents a linearized approximation of the interface cutting through the cell face. The $\alpha_f$ fractions are incorporated in the divergence operator, $\mathbf{D}$, as used in Eq.~\eqref{eq:pressure-poisson}.

For faces classified as ``solid'' (where the center is in the solid, $\psi < 0$) but possessing a non-zero fluid area fraction ($\alpha_f > 0$), we cannot solve for the velocity directly. Instead, we employ an extrapolated velocity value from a fluid neighbor to compute the flux. In edge cases where no suitable fluid neighbor exists for directional extrapolation (e.g., in under-resolved narrow gaps), we fall back to using the local boundary velocity $u_{bc}$ to avoid numerical instability and complex branching logic on the GPU.

For cells surrounded entirely by solid faces ($\alpha_f = 0$ for all faces), we interpolate values from the boundaries by solving a Laplace equation:
\begin{equation}
    \mathbf{D}_s \mathbf{G} p^{(k+1)} = 0
\end{equation}
Here, $\mathbf{D}_s$ denotes the divergence operator acting on the unweighted (intrinsic) velocities, distinguishing it from the weighted operator $\mathbf{D}$ extending to the cut-cells.

To formulate a unified equation for $\phi$ across both fluid and solid domains, we impose the following relation inside the solid:
\begin{equation}
    \delta \tilde{p}^{(k+1)} = c\,  \phi
\end{equation}
where $c$ is a constant of similar magnitude to the diagonal components of the matrices in Eq.~\eqref{eq:pressure_from_phi} ($c\approx\rho/\Delta t+ \rho v/\Delta x+\mu/\Delta x^2$). This yields:
\begin{equation}
    \mathbf{D}_s \mathbf{G} \, (p^{(k)} + c\,  \phi) = 0 \; \rightarrow \; \mathbf{D}_s \mathbf{G} \, \phi = - c^{-1} \, \mathbf{D}_s \mathbf{G} \, p^{(k)}
\end{equation}




\subsection{Staggered Grid Interpolation}
The convective fluxes require velocity components to be defined at locations where they are not naturally stored. To obtain the advective velocity $u$ at the face center, we employ a symmetric 4-point average. Similar to the flux calculation, if any node in this stencil lies within the solid, it is treated as a ghost value dependent on the local fluid velocity.

\begin{equation}
    \bar{u} = \frac{1}{4} \sum_{k=0}^{1} \sum_{j=0}^{1} u_{i+\delta_i j, j+\delta_j k, k} \quad \text{(Symmetric 4-point average)}
\end{equation}
where $u$ is the fluid value for valid cells, or the IBM-reconstructed value for solid cells (Eq.~\eqref{eq:ghost}). This ensures that the interpolated convective velocity field smoothly approaches the physical no-slip boundary condition ($u=0$), rather than artificially dropping to zero or using a staircase approximation at the solid interface.

\section{GPU Implementation Strategy}
\paragraph{Implementation Details}
Modern GPUs handle \texttt{min()}, \texttt{max()}, and floating-point comparisons more efficiently than warp-divergent \texttt{if} statements.
For the TVD scheme, the selection of the upwind node direction is determined branchlessly using \texttt{signbit} masks or comparison checks (e.g., $v>0$).
The Koren limiter evaluation involves computing the gradient ratio $r$ and clamping the result. This is implemented using intrinsic \texttt{fmin} and \texttt{fmax} instructions, which map directly to hardware units. For example, the limiter function can be evaluated as:
\begin{equation}
    \Psi(r) = \texttt{fmax}(0.0, \texttt{fmin}(2.0*r, \texttt{fmin}( (1.0+2.0*r)/3.0, 2.0 ) ) )
\end{equation}
To avoid division-by-zero when computing $r$, we use a protected division or multiply-add logic that gracefully handles uniform regions.
A key challenge in implementing Immersed Boundary Methods (IBM) on SIMT (Single Instruction, Multiple Threads) architectures is the potential for thread divergence caused by conditional branching (e.g., checking for cut-cells). To mitigate this, we adopt a \textit{bake-in} strategy: all geometric complexity is pre-calculated into sparse modification lists. The linear solver kernel itself remains agnostic to the boundary conditions, executing a uniform instruction set across the domain.

\subsection{Data Structures and Memory Layout}
To ensure coalesced memory access, we employ a Structure-of-Arrays (SoA) layout. The computational domain is discretized into $N_x \times N_y \times N_z$ cells. Consistent with visualization standards (e.g., VTK/ParaView), the $x$-index is the fastest-varying index in memory. 

The 7-point stencil coefficients for the linearized momentum equation are stored explicitly in seven aligned arrays:
\begin{equation}
    \mathcal{A}_{d} \in \mathbb{R}^{N_z \times N_y \times N_x}, \quad d \in \{C, W, E, S, N, B, T\}
\end{equation}
representing the coefficients for the Center, West, East, South, North, Bottom, and Top stencil points, respectively. The source term (Right-Hand Side) is stored in a separate array $\mathcal{B}$.

\subsection{Geometric Pre-computation}
To avoid conditional branching during the linear solve, we employ a unified "bake-in" strategy. During initialization, we identify all cells $c$ that require IBM corrections (i.e., cells with at least one solid neighbor). Due to the staggered grid arrangement, the geometric configuration differs for each velocity component. Therefore, we generate separate sparsity lists for each field: $\mathcal{L}_{IBM}^u, \mathcal{L}_{IBM}^v, \text{and } \mathcal{L}_{IBM}^w$.

For each cell in these lists, we pre-compute a set of scaling and modification factors for all three spatial directions ($d \in \{x, y, z\}$). Correct handling of the stencil requires modifying the center coefficient, the neighbor coefficients, and the source term.
For a given direction $d$ with neighbors at indices $nb,-$ and $nb,+$, we determine the status (e.g., Minus-Side Solid, Plus-Side Solid, Sandwiched, or None) and compute the following factors:
\begin{itemize}
    \item $K_{c,-}, K_{c,+}$: Factors for adding the neighbor's coefficient to the center.
    \item $M_-, M_+$: Self-multipliers for the neighbor coefficients.
    \item $X_-, X_+$: Cross-multipliers for adding one neighbor's original coefficient to the other.
    \item $B_{val,-}, B_{val,+}$: Terms for updating the source term $b_c$ with the inhomogeneous boundary values $u_{bc}$.
\end{itemize}
The factors $K, X,$ and $B_{val}$ implicitly contain the scaling ratio $\mathcal{R}_d = \mathcal{D}_{rescale}/\mathcal{D}_d$. For example, the cross-multiplier becomes $X \propto \mathcal{N}_{nb} \cdot \mathcal{R}_d$. A potential $0/0$ singularity arises if $\mathcal{D}_d \to 0$. However, by definition (Eq.~\ref{eq:D_rescale}), $|\mathcal{D}_{rescale}| \le |\mathcal{D}_d|$, ensuring that the numerator vanishes at least as fast as the denominator. In the numerical implementation, we compute this ratio using a safe division:
\begin{equation}
    \mathcal{R}_d = \begin{cases} 
    1 & \text{if } |\mathcal{D}_d| < \epsilon \\
    \mathcal{D}_{rescale} / \mathcal{D}_d & \text{otherwise}
    \end{cases}
\end{equation}
where $\epsilon$ is a small tolerance. This guarantees that the pre-computed factors remain bounded and valid even when the ghost-cell interaction point coincides with a grid node.

The global scaling factor $\mathcal{D}_{rescale}$ is also stored. Importantly, for directions that do not require modification (no solid neighbor), the additive factors ($K, X$) are set to 0 and the multiplicative factors ($M$) are set to $\mathcal{D}_{rescale}$, such that the mathematical operations remain valid but have no net effect beyond permissible scaling.

Complementary to the momentum operator, the discrete Laplacian for the pressure projection (Eq.~\ref{eq:pressure-poisson}) is also constructed in a "baked-in" fashion during the initialization phase. The face area fractions $\alpha_f$, required for the divergence and gradient operators, are computed from the SDF and incorporated directly into the 7-point pressure stencil coefficients. Since the solid boundaries are stationary, this pressure matrix is constant, avoiding the need for per-timestep geometric updates and allowing the use of highly optimized, uniform RB-GS kernels for the pressure solve.

\subsection{Time-Step Update Algorithm}
At each time step (or non-linear iteration), the linear system is assembled in three passes:

\paragraph{1. Base Stencil Generation}
A kernel launches over the full domain to compute the standard physical coefficients (convection, diffusion, time-accumulation) for a 7-point stencil, ignoring immersed boundaries.

\paragraph{2. Sparse Geometric Modification}
We launch a kernel over the list $\mathcal{L}_{IBM}$ to apply the geometric constraints. To ensure modularity and correct handling of the source term scaling, this is performed in two steps (or fused if the source term is available).

First, we modify the matrix coefficients $a$ and compute the inhomogeneous contribution $a_{inhom}$. For a cell $c$, we scale the center coefficient:
\begin{equation}
   a_c \leftarrow a_c \cdot \mathcal{D}_{rescale}
\end{equation}
Then, we loop over the three spatial directions $d \in \{x, y, z\}$ to update $a_c$ and the neighbors $a_{nb,\pm}$ using the factors $K, M, X$. Simultaneously, we accumulate the inhomogeneous term:
\begin{equation}
    a_{inhom,c} = \sum_{d, \pm} a_{nb,\pm} \cdot B_{val,\pm}
\end{equation}
where $B_{val}$ represents the term $\mathcal{N}_{bc} \mathcal{R} u_{bc}$.

Finally, in a separate pass (or at the end of the step), we apply the scaling to the source term vector $\mathbf{b}$ and subtract the inhomogeneous part to form the final right-hand side for the linear solver:
\begin{equation}
    b_c \leftarrow b_c \cdot \mathcal{D}_{rescale} - a_{inhom,c}
\end{equation}
This split ensures that the matrix operator can be constructed independently of the source term vector if needed (e.g., for implicit diffusion where the matrix is constant but the RHS varies).

\paragraph{3. Red-Black Gauss-Seidel Solver}
The linear solver kernel iterates over the domain using a Red-Black checkerboard pattern. Because the IBM logic has been fully absorbed into the modified coefficients $a$, the solver kernel contains no conditional branching related to geometry. It performs a standard update:
\begin{equation}
    u^{(k+1)}_c = \frac{1}{a_c} \left( b_c - \sum_{nb} a_{nb} \, u^{(k)}_{nb} \right)
\end{equation}
This approach ensures optimal thread occupancy and prevents warp divergence, as all threads execute identical floating-point operations.




\section{Verification}
To demonstrate the accuracy and robustness of the solver, we present a suite of verification benchmarks. These tests confirm the correct implementation of the numerical schemes and the immersed boundary treatment.

\subsection{Advecting Taylor-Green Vortex}
We verify the advection operator using the Taylor-Green vortex in a periodic cube domain $[0, 2\pi]^3$. The setup involves an initial vortex field superposed with a constant convective velocity $\vec{U} = (1, 1, 1)$. The pass criteria requires that the numerical kinetic energy decay matches the analytical profile $E(t)$, independent of the convection speed $\vec{U}$. This ensures the Galilean invariance and correctness of the advection scheme.

\subsection{Plane Poiseuille Flow}
We simulate a pressure-driven flow in a periodic domain containing a centered solid slab. This setup reproduces the classic plane Poiseuille flow between two parallel plates. We measure the resulting velocity profile against the analytical parabolic solution. It is crucial to note that the discrete velocities in our finite-volume formulation represent cell-averaged quantities. Therefore, to rigorously establish the second-order spatial accuracy ($O(\Delta x^2)$) of the discretization and IBM operations, we compare the numerical results against the cell-averaged analytical solution rather than point values.

\subsection{Flow Past Periodic Spheres}

To validate the solver's capability in handling complex porous media, we compute the flow through Simple Cubic (SC) arrays of spheres.  We used a periodic cubic domain of size $L$ and uniform grids with resolutions $N=16$, $32$, $64$, and $128$. The simulation was run for Stokes flow by setting $\rho=0$.

A signed-distance function is used to define the solid spheres inside (negative) and outside (positive) of the sphere with radius $R$. The solid volume fraction is varied by changing the radius of the sphere.
The solid volume fraction is given by $\phi = \frac{4/3 \pi R^3}{L^3}$, where $R$ is the radius of the spheres. The solid volume fraction is varied by changing the radius of the spheres.

On the fluid a body force, $f$, in the x-direction is applied to drive the flow. The superficial velocity, , $U_{sup}$, is computed by averaging the velocity field over the full domain (including the solid region where velocities equal 0). A dimensionless drag factor is defined as
\begin{equation}
    K = \frac{f \, L^3}{6 \pi \, \mu \, U_{sup} \, R}
\end{equation}
where $\mu$ is the dynamic viscosity.

Our results are compared with the reference values provided by Zick and Homsy (1982). These are numerical values, but highly accurate and widely used as a benchmark for immersed boundary methods.
Figure \ref{fig:drag_spheres} presents the comparison of our computed dimensionless drag against the reference data for solid volume fractions $\phi$ ranging from $0.001$ to $0.5236$. We observe excellent agreement across the entire range. Furthermore, we analyze the grid convergence of the method by computing the relative error with respect to a highly resolved reference case. As shown in Figure \ref{fig:drag_spheres} (b), the method exhibits clear second-order spatial accuracy ($O(N^{-2})$), confirming the efficacy of the proposed IBM treatment for curved geometries. For the highest resolutions the somewhat irregular behavior is due to the fact that the reference values are reported with a limited number of significant digits.

\begin{figure}[h!]
    \centering
    \begin{subfigure}[b]{0.48\textwidth}
        \centering
        \includegraphics[width=\textwidth]{figures/drag_simple_cubic_spheres.png}
        \caption{Dimensionless drag factor $F$ vs solid volume fraction $\phi$.}
        \label{fig:drag_spheres_plot}
    \end{subfigure}
    \hfill
    \begin{subfigure}[b]{0.48\textwidth}
        \centering
        \includegraphics[width=\textwidth]{figures/drag_simple_cubic_spheres_convergence.png}
        \caption{Relative error convergence showing $O(N^{-2})$ scaling.}
        \label{fig:drag_spheres_conv}
    \end{subfigure}
    \caption{Validation of flow through simple cubic arrays of spheres. (a) Comparison with Zick and Homsy (1982). (b) Grid refinement showing second-order accuracy.}
    \label{fig:drag_spheres}
\end{figure}

\subsection{Divergence and Energy Conservation}
Correct enforcement of the incompressibility constraint is monitored via automated tests. We track the $L_2$ norm of the velocity divergence field to ensure it remains within machine precision tolerances. Additionally, we verify the reduction of kinetic energy during the projection step, confirming the orthogonal projection of the velocity field onto the valid divergence-free subspace.

\end{document}

%%
%% End of file `elsarticle-template-1-num.tex'. 